Assume that in an environment we are interested in the probability of a random variable \rv[]{X}
conditioned on an observed evidence \rval[]{e}.
We have decided on a probabilistic model \M for the environment, and without any particular prior knowledge, we have
chosen a distribution $P$ from our model.

We ask an expert---a person familiar with the environment---to evaluate our distribution $P$.
Presumably, that expert has his own beliefs about the environment and thinks we should use a different
distribution $Q$, instead of $P$.
Consequently, when evaluating $P$, he would try to measure the loss of using the incorrect distribution $P(\XIe[])$,
while believing that the correct distribution is $Q$.

Let's assume that a real number can be used to measure this loss, and denote it as $\D[{\XIe[]}]$.
It seems reasonable as $P$ gets closer to $Q$ this measure decreases, and increases when $P$ diverges from $Q$.
Therefore, $\D[]$ effectively is a measure of the divergence of $P$ from $Q$, when $Q$ is the true distribution.

Interestingly, if we make certain assumptions about the properties of this divergence measure,
it turns out that---up to a constant factor---it must be equal to the Kullback–Leibler (KL) divergence.
Alternatively, by adopting a more minimal set of properties, we can prove that if we accept Shannon’s entropy as a
measure of the uncertainty of a distribution, we should also consider KL divergence as the appropriate method for
measuring divergence between distributions.

\begin{theorem}
    Suppose that $\D[]$ is a divergence measure, defined for probability distributions $Q$ and $P$,
    which satisfies the following properties:
    \begin{description}
        \item[Consistency with entropy] If $U$ is the uniform distribution with the same support
        as $Q$, then $\D[][Q][U] = \entropy{U}{} - \entropy{Q}{}$, where $\entropy{}{}$ denotes Shannon's
        entropy.
        \item[Additivity] For any two random variables $\XY$
        \[
            \D[\XY] = \D + \sumx \D[\YIx]Q(\rval[]{x})\cdot
        \]
    \end{description}
    Then, $\D[]$ is the Kullback-Leibler divergence.
\end{theorem}

\begin{proof}
    \newcommand{\T}{\varOmega}
    \renewcommand{\t}{\omega}
    \newcommand{\tIx}{\t \mid \rval[]{x}}
    \newcommand{\TIx}{\T \mid \rval[]{x}}
    \newcommand{\TX}{\T, \rv[]{x}}
    \newcommand{\tx}{\t, \rval[]{x}}
    \newcommand{\tent}[1]{\int q(#1) \ln q(#1) d\t}

    \newcommand{\Qx}{Q(\rval[]{x})}
    Let \rv[]{X} be a discrete random variable, and $P(\rv[]{X})$ and $Q(\rv[]{X})$ be two arbitrary distributions
    for \rv[]{X}.
    We define a new continuous random variable $\T$ and define
    \begin{equation*}
        p(\T = \t \mid \rvEq) =
        \begin{cases}
            \frac{1}{P(\rval[]{x})} & 0 \leq \t \leq P(\rvEq),\\
            0 & \text{otherwise}\cdot
        \end{cases}
    \end{equation*}
    Notice that both $P(\T \mid \rv[]{x})$ and $P(\TX)$ are uniform distributions, and
    $\entropy{P}{\TIx} = \ln P(\rval[]{x})$, $\entropy{P}{\TX} = 0$.

    We define $Q(\TIx)$ to be an arbitrary distribution that has the same support as $P(\TIx)$.
    This ensures that $Q(\TX)$ and $P(\TX)$ have the same support as well.
    Thus, we have
    \begin{align*}
        \Dcond[][][\TX][\TX] &= \entropy{P}{\TX} - \entropy{Q}{\TX}\\
        &= \sumx \tent{\tx}\\
        &= \sumx \Qx \int q(\tIx) \ln q(\tx) d\t,
    \end{align*}
    and
    \begin{align*}
        \Dcond[][][\TIx][\TIx] &= \entropy{P}{\TIx} - \entropy{Q}{\TIx}\\
        &= \ln P(\rval[]{x}) + \tent{\tIx}\cdot
    \end{align*}
    Now, we use the additivity property to obtain a formula for our divergence measure:
    \begin{align*}
        \D &= \D[\TX] - \sumx \D[\TIx]\Qx \\
        &= \sumx \Qx  \left( \int q(\tIx) \ln \frac{q(\tx)}{q(\tIx)}d\t - \ln P(\rval[]{x}) \right)\\
        &= \sumx \Qx  \left( \ln \Qx \int q(\tIx) d\t - \ln P(\rval[]{x}) \right)\\
        &= \sumx \Qx \ln \frac{\Qx}{P(\rval[]{x})}\cdot \qedhere
    \end{align*}
\end{proof}

In this theorem, the first property simply states that the entropy of a distribution is related to the divergence of
the uniform distribution.
The more the divergence is, the less the entropy should be.

The second property tells us that the divergence is an additive quantity.
We know that the divergence of $P(\XY)$ from $Q(\XY)$ is actually a measure of the generic loss of using $P(\XY)$
when we believe $Q(\XY)$ is the true distribution.
To calculate this loss, we can assume that \rv[]{X} happens before \rv[]{Y}.
In that way, this loss would include the loss of using $P(\rv[]{X})$, and then, if we observe
$\rv[]{X} = \rval[]{x}$, which can happen with probability $Q(\rval[]{x})$, it would also include the loss of
using $P(\YIx)$.
This property indicates that for combining these individual terms we should use addition.

Returning to our earlier discussion, this implies that to evaluate $P$ our expert should use the KL divergence
between $P$ and $Q$.
What if we find a distribution $C$ that has a small KL divergence from any other distribution in our model?
In that way, every expert would agree that our distribution is a good choice.
Interestingly, this relates to the concept of Bayesian convergence, where the posterior probability converges to the
same distribution regardless of the prior distribution.

\begin{definition}[$\epsilon$-convergence]
    In a probabilistic model \model, let \Qvar be a discrete random variable, and \Eset be a countable set of
    discrete random variables, all defined on \set{s}.
    We refer to \Qvar as the \emph{query variable}, and \Eset as the \emph{evidence set}.

    We say $\QvarIEset$ (\Qvar conditioned on \Eset) \emph{\econverges by $P \in \set{p}$},
    if for any random variable $\rv{e} \in \Eset$ and any probability measure $Q \in \set{p}$ we have:
    \begin{equation*}
        \Ex{Q}{\Dcond} < \epsilon,
    \end{equation*}
    where $\Ex{Q}{\,\cdot\,}$ denotes expectation relative to $Q$ and $\Dcond$ is the Kullback-Leibler divergence
    between $P(\XIe)$ and $Q(\XIe)$:
    \[
        \Dcond = \sum_{\rval{x}} Q(\xIe) \ln \frac{Q(\xIe)}{P(\xIe)}\cdot
    \]
\end{definition}

When \set{e} contains a single random variable, like \rv{e}, we simplify our notation.
Instead of writing $\Qvar|\{\rv{e}\}$ we simply use $\Qvar|\rv{e}$.

\begin{definition}[Divergence Measure]
    Let \Qvar be a query variable, and \Eset be an evidence set in a probabilistic model \model.
    We define the \emph{divergence measure of $P \in \set{P}$ for $\QvarIEset$} as follows:
    \begin{align}
        \label{eq:div-def}
        \DxP &= \maxQ \maxE \Ex{Q}{\Dcond} \\
        &= \maxQ \maxE \sumOverxe Q(\assign{x} , \assign{e}) \ln
        \frac{Q(\assign{x} \mid \assign{e})}{P(\assign{x} \mid \assign{e})} \cdot
    \end{align}
\end{definition}

Note that we usually simplify our notation and instead of writing the full assignment to a random variable, we only
write the assigned value.
For example, instead of $Q(\assign{x} , \assign{e})$ we write: $Q(\xe)$.

\begin{proposition}
    $\QvarIEset$ \econverges by $P$, if and only if:
    \begin{equation*}
        \DxP < \epsilon \cdot
    \end{equation*}
\end{proposition}

\begin{definition}[Central Distribution]
    In a model \model, we define the \emph{central distribution} for a query variable \Qvar, conditioned on the
    evidence set \Eset, to be a distribution $\central \in \set{P}$ which minimizes the divergence measure
    for $\QvarIEset$:
    \begin{equation*}
        \central = \arg\pmin \DxP\cdot
        \label{eq:center}
    \end{equation*}
\end{definition}

\begin{definition}[Aggregator]
    Consider two probabilistic models \model and \aggmodel.
    Let \Qvar be a query variable and $\Eset = \{\seqRvInf[e]\}$ be an evidence set in \M.
    Also, let \agg{\rv{e}}, \agg{\Qvar}, and \rv{i} be random variables defined on \agg{\set{s}} in \Magg.
    We refer to \agg{\rv{x}}, \agg{\rv{e}} as the \emph{aggregated} query and evidence variables, and \rv{i} as the
    \emph{selector} variable.

    Suppose that, for a probability measure $P \in \set{P}$, there is a probability measure
    $\agg{P} \in \agg{\set{P}}$, such that for any index \rval{i}, and every values \QvarVal, and \rval[i]{e}
    there is some \agg{\QvarVal}, and \agg{\rval{e}}, where we have
    \begin{equation*}
        P(\Qvar = \QvarVal, \rv[i]{e} = \rval[i]{e}) =
        \agg{P}(\agg{\Qvar} = \agg{\QvarVal}, \assign{\agg{e}} \mid \assign{I}) \cdot
        \label{eq:th1-1}
    \end{equation*}

    In addition, in \agg{P}, \agg{\Qvar} and \rv{I} are conditionally independent given \agg{\rv{E}}:
    \begin{equation*}
        \agg{P}(\agg{\Qvar} \mid \agg{\rv{E}}, \rv{I}) = \agg{P}(\agg{\Qvar} \mid \agg{\rv{E}}) \cdot
    \end{equation*}

    We call \agg{P} an \emph{aggregator of $P$ with respect to $\QvarIEset$}.

\end{definition}

\begin{definition}[Aggregate Model]
    \label{def:aggregate-model}
    Consider a probabilistic model \model.
    Let \Qvar be a query variable, and \Eset be an evidence set in \M.
    We say a model \aggmodel is the \emph{aggregate model} of \M with respect to $\QvarIEset$,
    If every probability measure in \Magg is the aggregator of a probability measure in \M with respect to $\QvarIEset$.
    %, and every probability measure in \M has at least one aggregator in \Magg with respect to $\QvarIEset$.
\end{definition}

\begin{definition}[Selective Aggregate Model]
    Let \Qvar be a query variable, and \Eset be an evidence set in \model.
    We say that \Magg, an aggregate model of \M with respect to $\QvarIEset$, is \emph{selective}, if for
    every $P \in \set{P}$ and any value \rval{i} of the selector variable \rv{I}, there is an aggregator of
    $P$ in \Magg, such that $\agg{P}(i) = 1$.

\end{definition}


\begin{theorem}
    Let \Qvar be a query variable, and \Eset be an evidence set in a probabilistic model \model.
    Suppose that \aggmodel is a selective aggregate model of \M with respect to $\QvarIEset$.
    Let \agg{\rv{x}}, \agg{\rv{e}} be the aggregated query and evidence variables in \Magg.
    If $\agg{P} \in \agg{\set{P}}$ is an aggregator of the probability measure $P \in \set{P}$,
    with respect to $\QvarIEset$, then we have
    \begin{equation*}
        \DxP = \aggDxP \cdot
    \end{equation*}

    In addition, the central distribution for $\agg{\Qvar}|\agg{\rv{e}}$ is an aggregator of the central
    distribution for $\QvarIEset$.
\end{theorem}

\begin{proof}
    \newcommand{\aggExpKL}[1]{\ExpKL[\agg{\rval{x}}][\agg{\rval{e}}][\agg{#1}][\agg{P}]}
    \newcommand{\aggxe}{\agg{\QvarVal}, \agg{\rval{e}}}
    \newcommand{\aggxIe}{\agg{\QvarVal} \mid \agg{\rval{e}}}
    \newcommand{\aggXIE}{\agg{\Qvar} \mid \agg{\rv{e}}}
    Let
    \begin{equation*}
        \agg{R} = \arg\max_{\agg{Q} \in \agg{\set{P}}} \aggExpKL{Q}.
    \end{equation*}
    Since \Magg is an aggregate model, $\agg{R}$ must be an aggregator of some probability measure $R \in \set{P}$.
    It follows that
    \begin{align*}
        \aggDxP
        &= \aggExpKL{R} \\
        &= \sum_i \agg{R}(\rval{i}) \ExpKL[][{\rval[i]{e}}][R] \\
        &\leq \max_i \ExpKL[][{\rval[i]{e}}][R] \cdot
    \end{align*}
    This implies
    \begin{equation}
        \aggDxP \leq \DxP \cdot
        \label{eq:th-2-proof-1}
    \end{equation}

    Similarly, let
    \begin{equation*}
    (R,k)
        = \arg\max_{(Q,i)} \ExpKL[][{\rval[i]{e}}] \cdot
    \end{equation*}
    \Magg is a selective aggregate model.
    This means that $R$ must have an aggregator \agg{R} in \Magg, which selects $k$ by assigning it a
    probability of 1: $\agg{R}(k) = 1$.
    Therefore,
    \begin{align*}
        \DxP
        &= \ExpKL[][{\rval[k]{e}}][R] \\
        &= \sum_{\aggxe} \agg{R}(\aggxe \mid k) \ln\frac{\agg{R}(\aggxIe, k)}{\agg{P}(\aggxIe, k)} \\
        &= \aggExpKL{R} \\
        &\leq \aggDxP \cdot
    \end{align*}
    By Equation~\ref{eq:th-2-proof-1}, we conclude that
    \begin{equation*}
        \DxP = \aggDxP \cdot
    \end{equation*}

    To prove the second part of the theorem, let \central[\aggXIE] be a central distribution
    for $\agg{\Qvar}|\agg{\rv{e}}$.
    Since \Magg is an aggregate model for \M, \central[\aggXIE] must be an aggregator for a probability
    measure $\central \in \agg{\set{P}}$.
    If we assume, by contradiction, that \central is not a central distribution for $\QvarIEset$ in \M, then there must
    be a probability
    measure $R \in \set{P}$ such that:
    \begin{equation*}
        \DxP[R] < \DxP[\central] = \aggDxP[C_{\aggXIE}] \cdot
    \end{equation*}
    On the other hand, since \Magg is a selective model, $R$ must have an aggregator $\agg{R} \in \agg{\set{P}}$,
    and we have
    \begin{equation*}
        \DxP[R] = \aggDxP[R] > \aggDxP[{C_{\aggXIE}}],
    \end{equation*}
    which is a contradiction.
\end{proof}

\begin{proposition}
    The central distribution for $\agg{\Qvar}|\agg{\rv{e}}$ is an aggregator of the central distribution
    for $\QvarIEset$.
\end{proposition}

\begin{theorem}[Aggregate Model]
    Consider two parametric probabilistic models \pmodel and \aggmodel.
    Let \Qvar be a query variable and $\Eset = \{\seqRvInf[e]\}$ be an evidence set in \Mt.
    Also, let \agg{\rv{e}}, \agg{\Qvar}, \aggParam, and \rv{i} be random variables defined
    on \agg{\set{s}} in \Magg.

    Suppose that for any \rval{i}, and every values \QvarVal, \rval[i]{e}, and \paramValue, there is some
    \agg{\QvarVal}, \agg{\rval{e}}, and \aggParamVal, such that we have
    \begin{equation*}
        \pi(\Qvar = \QvarVal, \rv[i]{e} = \rval[i]{e}  \mid \param = \paramValue) =
        \agg{\pi}(\agg{\Qvar} = \agg{\QvarVal}, \assign{\agg{e}}
        \mid \aggParam = \aggParamVal, \assign{I}),
        \label{eq:th1-1}
    \end{equation*}
    where $\pi$ is the likelihood distribution in \Mt, and $\agg{\pi}$ is the likelihood distribution in \Magg.

    In addition, in every probability measure $\agg{P} \in \agg{\set{P}}$, \aggParam and $\rv{i}$ are independent:
    \begin{equation*}
        \agg{p}(\aggParam \mid \rv{I}) = \agg{p}(\aggParam),
    \end{equation*}
    while the prior distribution $\agg{p}(\aggParam)$, is the same as the prior distribution
    of \param induced by some probability measure $P \in \set{P}$:
    \begin{equation*}
        p(\param = \paramValue) = \agg{p}(\aggParam = \aggParamVal)\cdot
    \end{equation*}
    Also, \agg{\Qvar} and \rv{I} are conditionally independent given \agg{\rv{E}}:
    \begin{equation*}
        \agg{P}(\agg{\Qvar} \mid \agg{\rv{E}}, \rv{I}) = \agg{P}(\agg{\Qvar} \mid \agg{\rv{E}})
    \end{equation*}

    On the other hand, for every probability measure $P \in \set{P}$, and any arbitrary distribution $Q$,
    there is a $\agg{P} \in \agg{\set{P}}$ such that for every \rval{i}, and \paramValue there exists a \aggParamVal
    such that
    \begin{equation*}
        \agg{p}(\aggParamVal, \rval{I}) = p(\paramValue)Q(\rval{I})\cdot
    \end{equation*}
    We refer to \Magg as the aggregate model of \Mt for $\QvarIEset$.

    If these conditions hold, then for any $P \in \set{P}$, there is a $\agg{P} \in \agg{\set{P}}$ such that we have
    \begin{equation*}
        \DxP = \aggDxP,
    \end{equation*}
    and vice versa.
\end{theorem}
\begin{proof}
    sdfsfs
\end{proof}
\begin{proposition}
    sdfs
\end{proposition}
\begin{theorem}

    If these conditions hold, then for any $P \in \set{P}$ and $\agg{P} \in \agg{\set{P}}$ we have
    \begin{equation*}
        \DxP = \measure{D}_{\agg{\Qvar} \mid \agg{\rv{e}}} \left[ \agg{P} \right] \cdot
    \end{equation*}
    We refer to \Magg as the aggregate of \Mt for $\QvarIEset$.
    \begin{equation*}
        \agg{P}(\agg{\Qvar} \mid \agg{\rv{e}}) = \agg{P}(\agg{\Qvar} \mid \agg{\rv{e}}, \rv{i})\cdot
        \label{eq:th1-2}
    \end{equation*}
    Then for every $P \in \set{P}$:
    \begin{equation*}
        \DxP = \measure{D}_{\agg{\Qvar} \mid \{\agg{\rv{e}} \}} \left[ \agg{P} \right] \cdot
    \end{equation*}
\end{theorem}

\begin{proof}
    Let $Q_m$ be the distribution, and $\rv[m]{e}$ be the evidence that maximizes $\Ex{Q}{\Dcond}$.
    Therefore, by Equation~\ref{eq:div-def} we have
    \begin{equation*}
        \DxP = \ExpKL[][e_m][Q_m]\cdot
    \end{equation*}

    The definition of $\DxP$ involves
    maximizing with respect to
    Assume, by contradiction, that $\given{S}{\Eset}$ does not $\epsilon$-converge by $P$ in $\M$.
    Therefore, there must exist $Q$ and $E_i \in \Eset$ in $\M$ such that:


\end{proof}


\begin{definition}[Conditional Divergence Measure]
    \label{def:cond_divergence}
    Assume that \pmodel is a parametric model.
    Let \Qvar be a query variable and \Eset be an evidence set in \Mt.
    We define the \emph{divergence measure of $P \in \set{P}$ given $\paramValue$ for $\QvarIEset$} as follows:
    \begin{align}
        \DxPcond &= \maxE \DExparam \\
        &= \maxE \Dparam \cdot
        \label{eq:cond_div_measure}
    \end{align}
\end{definition}

In this definition, the query set $\Xpairs$ plays a crucial role.
Specifically, \set{x} defines the set of random variables for which we are interested in their distributions.
Each \Qvar represents a random variable whose distribution is important to us, and each \rv{e} corresponds to the
evidence that is available for that variable.

For a parametric model (as defined in Definition~\ref{def:param-model}), it seems reasonable to define a conditional
form of our divergence measure, conditioned on the value of the parameter.

To do so, in Equation~\ref{eq:div_measure}, we replace normal expectation with conditional expectation.
There will be no longer a need to maximize over different distributions, since in a parametric model the parameter
fully determines the distribution of the sample space.

Even if we cannot find a distribution that has a small divergence from every distribution, still using the distribution
that minimizes the maximum divergence seems a reasonable choice.

In a parametric model, the maximum of divergence happens for a prior distribution
that is concentrated at a single value of the parameter $\paramValue$.
That means, for a parametric model which allows any high peak prior, instead of maximizing over the set of
all acceptable distributions, we can simply maximize over the parameter space, using the conditional version
of the divergence measure.

\begin{lemma}
    \label{lm:DxP_uppper_bound}
    Assume that \pmodel is a parametric model.
    Let \Qvar be a query variable and \Eset be an evidence set in \Mt.
    For every $P \in \set{P}$ we have
    \begin{equation*}
        \DxP \leq \tmax \DxPcond
    \end{equation*}
\end{lemma}

\begin{proof}
    Since the Kullback-Leibler divergence is positive, its expected value is also positive, and for any arbitrary
    $R \in \set{P}$, $\rv{E} \in \Eset$, and parameter value $\paramValue$ we have
    \[
        \DExparam[][R] = \Dparam[R] \geq 0\ \cdot
    \]
    Using this, for every distribution $P$, we obtain
    \newcommand{\RP}{\sumOverxe \pi(\xeIt) \ln \frac{R(\xIe)}{P(\xIe)}}
    \[
        \RP \leq \Dparam \cdot
    \]
    Since this inequality holds for any $\paramValue$ and any $\rv{e} \in \Eset$, the maximum of the right-hand side
    w.r.t. $\paramValue$ and $\rv{e}$ is also greater than or equal to the maximum of the left-hand side:
    \begin{equation}
        \tmax \maxE \RP \leq \tmax \DxPcond \cdot
        \label{eq:ineq-theorem1}
    \end{equation}

    \Mt is a parametric model and $\paramDist[R]{\xe}$, where $r(\paramValue)$ is
    a prior probability distribution.
    Trivially, we have
    \[
        \RP \leq \tmax \RP\cdot
    \]
    This inequality holds for every value of \paramValue.
    Therefore, we can multiply this inequality by $r(\paramValue)$, which is non-negative, and integrate.
    Notice that the right hand side is not a function of $\paramValue$ and it remains unchanged
    because $\int \dt[r] = 1$:
    \[
        \Ex{R}{\Dcond[R]} \leq \tmax \RP\cdot
    \]
    For any random variable $\rv{e} \in \Eset$, this inequality holds.
    That implies
    \[
        \maxE \Ex{R}{\Dcond[R]} \leq \maxE \tmax \RP\cdot
    \]
    Now from Inequality~\ref{eq:ineq-theorem1}, for \emph{any} distribution $P$ and $R$ in \Mt, we have
    \[
        \Ex{R}{\Dcond[R]} \leq \tmax \DxPcond,
    \]
    which proves the lemma.
\end{proof}

\begin{theorem}
    \label{th:Dx_parametric_form}
    Suppose \pmodel is a parametric model.
    Let \Qvar be a query variable, and \Eset be an evidence set in \Mt.
    If for every value of $\paramValue$ there is a distribution $R \in \set{P}$ such that
    $R(\rval[]{s}) = \pi(\rval[]{s} \mid \paramValue)$ for every $\rval[]{s} \in \set{S}$,
    then for any $P \in \set{P}$ we have:
    \begin{equation}
        \DxP = \tmax \DxPcond\cdot
        \label{eq:div_parametric}
    \end{equation}
    \label{th:div_for_parametric}
\end{theorem}

\begin{proof}
    For an arbitrary distribution $P \in \set{P}$, we let
    \[
        \paramValue^* = \arg\tmax \DxPcond \cdot
    \]
    In \Mt there should be a distribution $R \in \set{P}$ that equals $\pi(\rval[]{s} \mid \paramValue^*)$
    for every $\rval[]{s} \in \set{S}$, and therefore it defines the same distribution for any random variable.
    It follows that for every $\rv{e} \in \Eset$, and every $P \in \set{p}$
    \begin{equation}
        \label{eq:th1}
        \Ex{R}{\Dcond[R]} = \DExparam[\paramValue^*] \cdot
    \end{equation}
    On the other hand, clearly
    \[
        \maxQ \maxE \Ex{Q}{\Dcond} \geq \maxE \Ex{R}{\Dcond[R]}\cdot
    \]
    Notice that the left hand side is $\DxP$, and by Equation~\ref{eq:th1}
    \[
        \DxP \geq \maxE \tmax \DExparam\cdot
    \]
    This inequality and Lemma~\ref{lm:DxP_uppper_bound} prove the theorem.
\end{proof}

\begin{definition}[Generic Parametric Model]
    We say a parametric model \pmodel is \emph{generic}, if for any arbitrary probability distribution $r$,
    there is a probability measure $R \in \set{P}$, such that for every $\rval[]{s} \in \set{s}$
    \[
        \paramDist[r]{\rval[]{s}}\cdot
    \]
\end{definition}

\begin{proposition}
    A generic parametric model satisfies the conditions of Theorem~\ref{th:Dx_parametric_form}.
\end{proposition}

\begin{theorem}[Minimax theorem~\cite{Du1995}]
    \label{th:minimax}
    Let $\set{X} \subset \mathbb {R}^{n}$ and $\set{Y} \subset \mathbb {R}^{m}$ be compact convex sets.
    If $f:\set{X}\times \set{Y}\rightarrow \mathbb {R}$ is a continuous function that is concave-convex, i.e.

    \begin{itemize}
        \item $f(\cdot ,y):\set{X}\to \mathbb {R}$ is concave for fixed $y$, and
        \item $f(x,\cdot ):\set{Y}\to \mathbb {R}$ is convex for fixed $x$.
    \end{itemize}
    Then we have that
    \begin{equation}
        \max_{x\in \set{X}}\min_{y\in \set{Y}}f(x, y) = \min_{y\in \set{Y}}\max_{x\in \set{X}}f(x,y)\cdot
        \label{eq:minimax}
    \end{equation}
\end{theorem}

\begin{theorem}
    \label{th:central_by_max}
    Consider a generic parametric model \pmodel,
    and let $\Xpairs$ be a set of pairs of discrete random variables in \Mt.
    Assume that for every $Q_1, Q_2 \in \set{P}$ and any real number $0 \leq \alpha \leq 1$, there is
    a distribution $Q^* \in \set{P}$, such that for every $i, \xe$ we have
    \begin{equation}
        Q^*(\xIe) = \convexCond{\xIe}\cdot
        \label{eq:convex_dist_set}
    \end{equation}
    In this model, the central distribution for $\set{X}$ can be obtained by:
    \begin{align}
        \central &= \arg\maxP \int \DxPcond \dt\\
        &= \arg\maxP \Dminimax,
        \label{eq:Cx_by_max}
    \end{align}
    where $\paramDist[p]{\rval[]{s}}$, for every $\rval[]{s} \in \set{s}$.
\end{theorem}

\begin{proof}
    \newcommandx{\f}[2][1 = r, 2 = P_{\XIE}]{ f(#1, #2)}
    We prove this theorem using the Minimax theorem.
    For that, we define a functional as follows:
    \begin{equation*}
        \f = \int \DExparam \dt[r],
    \end{equation*}
    where $r(\paramValue)$ is an arbitrary probability distribution, and $P \in \set{P}$.
    Since $r(\paramValue)$ can be any distribution, maximizing with respect to $r$ is similar to maximizing
    with respect to $\paramValue$, and we have
    \[
        \central = \arg\pmin \max_{r} \f\cdot
    \]

    The domain of $r$ is the set of all probability distributions, which is a compact convex set.
    Moreover, from Equation~\ref{eq:convex_dist_set}, it follows that the set of probability distributions for
    $\Qvar|\rv{E}$ in \Mt is also a compact convex set.
    Thus, the domain of $\f$ is a compact convex set.

    For a fixed $P_{\XIE}$, $\f[\cdot]$ is a linear function, which is both concave and convex.
    On the other hand, if we let $\paramDist[r]{\xe}$ for $\f[r][\cdot]$ we have
    \begin{equation}
        \f[r][\cdot] = c - \sumOverixe R(\xe) \ln P(\xIe)\cdot
        \label{eq:max_for_P}
    \end{equation}
    If we calculate the Hessian matrix of $\f[r][\cdot]$, we observe that the matrix is diagonal, and each element
    of the main diagonal is $\frac{R(\xIe)}{P(\xIe)^2}$.
    Thus, the hessian matrix of $\f[\cdot]$ is positive semi-definite and $\f[r][\cdot]$ is a convex function.
    Now, Theorem~\ref{th:minimax} implies
    \begin{equation}
        \pmin \max_{r} \f = \max_{r} \pmin \f\cdot
        \label{eq:divergence_minmax}
    \end{equation}
    From Equation~\ref{eq:max_for_P} we can see that for an arbitrary fixed $r$ we have
    \[
        \pmin \f = \maxP \sumOverixe R(\xe) \ln P(\xIe)\cdot
    \]

    We know that the Kullback–Leibler divergence, and therefore its expected value, is non-negative:
    \[
        \Ex{R}{\sum_i \Dcond[R]} \geq 0\cdot
    \]
    By simple algebra we get
    \[
        \sumOverixe R(\xe) \ln P(\xIe) \leq \sumOverixe R(\xe) \ln R(\xIe)\cdot
    \]
    Recall that \Mt is a generic parametric model, so $R \in \set{P}$, and we can reach this
    upper bound by letting $P = R$:
    \[
        \pmin \f = \Dminimax[R]\cdot
    \]
    By Equation~\ref{eq:divergence_minmax} we conclude
    \[
        \central = \arg\max_{R \in \set{R}} \int \DExparam[][R] \dt[r]\cdot \qedhere
    \]
\end{proof}

\begin{theorem}
    \label{th:lambda}
    Under the assumptions of Theorem~\ref{th:central_by_max}, the central distribution for \set{X} satisfies:
    \begin{equation}
        \DxPcond[\central] = \lambda,
        \label{eq:lambda}
    \end{equation}
    where $\lambda$ is a constant.
\end{theorem}
\begin{proof}
    By Theorem~\ref{th:central_by_max}, the central distribution for \set{X} maximizes
    \[
        \functional{p(\paramValue)} =\Dminimax,
    \]
    with respect to different choices of the prior distribution $p(\paramValue)$, subject to the
    constraint $\int \dt[p] = 1$.

    Applying the Lagrange multiplier technique, the local maximum of $\functional{}$ is a stationary point of the
    Lagrangian functional:
    \[
        \Lagr[p(\paramValue)] = \Dminimax - \lambda \left( \int \dt - 1 \right)\cdot
    \]
    Let us denote the prior distribution corresponding to the central distribution \central by $\cPrior$.
    Let $p(\paramValue) = \cPrior + \epsilon \eta(\paramValue)$ be a variation of $\cPrior$, where $\eta(\paramValue)$ is
    a continuously differentiable function with $\eta(0) = \eta(1) = 0$, and $\epsilon$ is a small constant.

    Since $\cPrior$ is an extremizing function for $\functional{p(\paramValue)}$, the lagrangian
    $\Lagr$ has a stationary point where $\epsilon = 0$, and we have
    \[
        \parDerivAt[\Lagr] = 0\cdot
    \]
    Note that
    \begin{align*}
        \de{p(\paramValue)} &= \eta(\paramValue),\\
        \de{P(\xe)} &= \margDist[\eta]{\xe},\\
        \de{P(\rval{e})} &= \margDist[\eta]{\rval{e}}\cdot
    \end{align*}
    Using the chain rule, we calculate the partial derivative of $\Lagr$ relative to $\epsilon$:
    \begin{align*}
        \parDeriv\ =\ &\int \left( \Dparam  \right) \dt[\eta] \\
        &- \int \left( \sumOverixe \pi(\xeIt) \frac{\margDist[\eta]{\xe}}{P(\xe)} \right) \dt \\
        &+ \int \left( \sum_i \sum_{\rval{e}} \pi(\eIt) \frac{\margDist[\eta]{\rval{e}}}{P(\rval{e})} \right) \dt \\
        &- \int \lambda \dt[\eta]
    \end{align*}
    Since $\paramDist{\xe}$ and $\paramDist{\rval{e}}$ we get
    \begin{equation*}
        \parDeriv = \int \left(-\lambda + \Dparam \right) \dt[\eta]\cdot
    \end{equation*}
    We let $ \parDerivAt = 0$.
    Since $\eta(\paramValue)$ is an arbitrary function, by applying the fundamental lemma of the calculus of
    variations and letting $\epsilon = 0$, for every value of $\paramValue$ we have
    \[
        \Dparam[\central] = \lambda\cdot \qedhere
    \]
\end{proof}

Consider a scenario where we are trying to predict the outcome of a four-sided dice. This dice is unfair, with each
face marked by a binary number: $00$, $01$, $10$, and $11$.
Before making our prediction, we can roll the dice for several times and use the results to learn the associated
probabilities.
We assume that the dice rolls are independent.

Let's focus on determining the central distribution for this problem.
The first step in this process is to establish our query set $\Xpairs$.

In a prediction problem, we are interested in the distribution of the random variable we want to predict,
conditioned on the observed data.
For example, if we roll the dice for three times before making our prediction,
we will need $P(\rv[4]{x} \mid \rv[1]{x},\rv[2]{x},\rv[3]{x})$, where \Qvar denotes the outcome
of the $i$-th dice roll.

Thus, one natural way for defining a query set $\Xpairs$ is to let each \Qvar be the outcome of
the $i$-th dice roll, and define each \rv{e} as a vector of the previous dice outcomes $(\seqRv{i-1})$.

Now let us calculate the conditional divergence measure for \set{x} as defined in Definition~\ref{def:cond_divergence}:
\newcommand{\seqvalxx}[1][x]{\rval[1]{#1}, \rval[2]{#1}}
\begin{align*}
    \DxPcond &= \KLcond{\rval[1]{x}}{} \\
    &+ \KLcond{\rval[2]{x}}{\rval[1]{x}} \\
    &+ \KLcond{\rval[3]{x}}{\seqvalxx} \\
    &+ \ \cdots \\
    &+ \KLcond{\rval[n]{x}}{\seqRval{n-1}}\cdot
\end{align*}
After factoring out $\pi(\seqRval{n} \mid \paramValue)$ from every term, we use the chain rule for conditional
probability:
\begin{align*}
    \DxPcond &=
    \sum_{\seqRval{n}} \pi(\seqRval{n} \mid \paramValue) \ln
    \frac
    {\pi(\rval[1]{x} \mid \paramValue) \pi(\rval[2]{x} \mid \rval[1]{x}, \paramValue)
    \cdots \pi(\rval[n]{x} \mid \seqRval{n-1}, \paramValue)}
    {P(\rval[1]{x}) P(\rval[2]{x} \mid \rval[1]{x}) \cdots P(\rval[n]{x} \mid \seqRval{n-1})} \\
    &= \KLcond{\seqRval{n}}{} \cdot
\end{align*}
By Equation~\ref{eq:Cx_by_max}, for the central distribution we have
\begin{align*}
    \central &= \arg\max_{P} \int \left(\KLcond{\seqRval{n}}{}\right) \dt \\
    &= \arg\max_{P} \int \sum_{\seqRval{n}} p(\seqRval{n}, \paramValue) \ln
    \frac{p(\paramValue \mid \seqRval{n})}{p(\paramValue)} d\paramValue \cdot
\end{align*}

That implies, the central distribution for our problem is the same as the distribution we obtain when using the
reference prior~\cite{BERNARDO200517} for our main parameter.

Imagine playing a game of chance involving our four-sided dice.
After rolling the dice, we input the outcome into a mysterious machine.
The machine then announces whether we have won or lost the game.
We have no clue how the machine works. We are not even sure if the dice outcome truly affects the machine’s output.
All we know is that the machine is memory-less, and the result of each game is independent.

Our goal is to learn how the machine operates through experimentation.
We aim to be able to predict whether we will win the game after observing the dice outcome.

We assume that, first, we roll the dice $n$ times, and record each outcome.
Next, we input some of these dice outcomes into the machine and observe the result of each game.
As we will see, this assumption is important for deriving the central distribution more easily.

Let us define a query set for this problem.
We denote this set as $\Xpairs$, where each \Qvar is the result of the $i$-th game, and the evidence \rv{e} is a
vector
\[
    (\seqRv{i-1}, \seqRv[y]{n}),
\]
where each \rv{y} denotes the $i$-th dice outcome.

The conditional divergence for \set{x} can be obtained by
\begin{align*}
    \DxPcond &= \KLcond{\rval[1]{x}}{\seqRval[y]{n}} \\
    &+ \KLcond{\rval[2]{x}}{\rval[1]{x},\seqRval[y]{n}} \\
    &+ \KLcond{\rval[3]{x}}{\seqvalxx, \seqRval[y]{n}} \\
    &+ \ \cdots \\
\end{align*}
After using the chain rule for conditional probability the right hand side simplifies to
\begin{align*}
    \KLcond{\seqRval{n}}{\seqRval[y]{n}} \cdot
\end{align*}

Consequently, an alternative way for defining our query set is to define it as a single
pair $\left\{(\rvec{x}, \rvec{e})\right\}$, where \rvec{x} is the vector of $n$ game's result, and \rvec{e} is the vector of $n$
dice outcome.

Using the new definition for \set{x}, by Equation~\ref{eq:lambda} we have:
\begin{equation*}
    \KLcond[\central]{\vval{x}}{\vval{e}} = \lambda
\end{equation*}
\emph{derive a $\paramValue = (\paramA, \paramB)$ closed form for central dist...}

\[
    \frac{\pi(\vval{x} \mid \vval{e}, \paramValue)}{\central(\vval{x} \mid \vval{e})}
    = \frac{\pi(\paramValue \mid \vval{x}, \vval{e})}{\central(\paramValue \mid \vval{e})}
\]

sdfd

\[
    \KLabVec
\]

Assume that we decide to simplify our learning task.
We choose to consider only the least significant digit of the dice outcome and ignore the other digit.
For example, if the dice outcome is $10$, we treat it as $0$.
This approach decreases the number of possible outcomes of the dice, making the learning process simpler.
However, it also reduces our ability to predict the game’s result accurately.

In order to find a central distribution for this version of our game, we should slightly modify the definition of the
query set.
While each \Qvar remains unchanged, we need to represent each \rv{e} as a vector
\[
    (\seqRv{i-1}, \seqRv[y']{n}),
\]
where each \rv{y'} represents the least significant digit of the dice outcome.
We denote this new query set as \set{x'} and the previous query set as \set{x}.

Following similar steps as we did for \set{x}, we can show that \set{x'} also has a more compact counterpart,
which includes a single pair $(\rvec{x}, \rvec{e'})$, where \rvec{x} is the vector $(\seqRv{n})$, and \rvec{e'} is
the vector
$(\seqRv[y']{n})$.

\emph{explain how the central distribution for this version is different...}

When we have a lot of learning data, there is no point in using this simpler version of our learning task.
However, when the available data is not enough for the more complex version, the simpler version can still produce
relatively meaningful results,
while the data is too small for the more complex version.

Is it possible to have a probabilistic inference framework that is able to automatically choose between these two
versions of the problem based on the quality and amount of available data?
It seems that the central distributing has this ability.
For that, we need to use the central distribution for the query set $\set{x} \cup \set{x'}$.
This seems reasonable, since when we want our probabilistic inference framework to choose between two versions of
our problem, our query set should include random variables of both versions.

\begin{theorem}
    Let \pmodel be a generic parametric model, and $\Xsingles$ be a set of random variables in \Mt.
    Let $\paramFunc$, for every $i$, be a function of the parameter $\param$ that fully parametrizes \Qvar,
    so that for $\paramFuncValue = \paramFunc(\paramValue)$ we have:
    \begin{equation}
        \pi(\xIt) = \pi \left( \rval{x} \mid \paramFuncValue \right)\cdot
        \label{eq:func_of_theta}
    \end{equation}
    The prior distribution corresponding to the central distribution for \set{X} in \Mt can be characterized
    by the following fixed point equation:
    \begin{equation}
        \cPrior \propto \left( \prod_i \cMass{\tIw} \rhoPrior \right)^{\frac{1}{|\set{x}|}},
        \label{eq:fixed_point_cx}
    \end{equation}
    where
    \[
        \rhoPrior \propto \exp \left\{\sum_{\rval{x}} \pi(\xIw) \ln \cMass{\wIx} \right\}\cdot
    \]

\end{theorem}

\begin{proof}
    First, we show that the set of distributions of each \Qvar is a convex set.
    Let $Q^*(\rval{x}) = \convexCond{\rval{x}}$, for some $Q_1, Q_2 \in \set{P}$, $0 \leq \alpha \leq 1$ and
    arbitrary $i, \rval{x}$.

    \Mt is a parametric model.
    Therefore,
    \begin{align*}
        Q^*(\rval{x}) &= \alpha \margDist[q_1]{\rval{x}} + (1 - \alpha) \margDist[q_2]{\rval{x}}\\
        &= \int \pi(\xIt) (\convexCond[q]{\paramValue}) d\paramValue\cdot
    \end{align*}
    $\convexCond[q]{\paramValue}$ is a valid prior distribution and \Mt is a generic parametric model.
    That implies $Q^* \in \set{P}$.

    Thus, by Theorem~\ref{th:lambda} the central distribution of \set{X} satisfies
    \begin{equation}
        \sumOverix \pi(\xIt) \ln \frac{\pi(\xIt)}{\central(\rval{x})} = \lambda\cdot
        \label{eq:lambda_th_proof}
    \end{equation}
    Moreover, from Equation~\ref{eq:func_of_theta} and Bayes' theorem it follows
    \begin{align*}
        \sumOverix \pi(\xIt) \ln \frac{\pi(\xIt)}{\central(\rval{x})}
        &= \sumOverix \pi(\xIw) \ln \frac{\pi(\xIw)}{\central(\rval{x})}\\
        &= \sumOverix \pi(\xIw) \ln \frac{\cMass{\wIx}}{\cPrior[(\paramFuncValue)]}\\
        &= \sum_i \ln \frac{\rhoPrior}{\cPrior[(\paramFuncValue])}\cdot
    \end{align*}
    Since each $\paramFunc$ is a function of $\param$, for $\paramFunc(\paramValue) = \paramFuncValue$
    we have $\cMass{\paramValue, \paramFuncValue}=\cPrior$.
    Therefore,
    \begin{equation*}
        \cPrior[(\paramFuncValue)] = \frac{\cPrior}{\cMass{\tIw}}\cdot
    \end{equation*}
    Now we rewrite Equation~\ref{eq:lambda_th_proof}:
    \begin{align*}
        |\set{X}| \ln \cPrior &= \sum_i \ln \left(\cMass{\tIw} \rhoPrior \right) - \lambda\\
        \cPrior &= \left( e^{-\lambda} \prod_i \cMass{\tIw} \rhoPrior \right)^{\frac{1}{|\set{x}|}} \qedhere
    \end{align*}
\end{proof}


\begin{theorem}
    Let \pmodel be a generic parametric model, and $\Xpairs$ be a set of random variables in \Mt.
    Let $\paramFunc$, for every $i$, be a function of the parameter $\param$ that fully parametrizes \Qvar,
    so that for $\paramFuncValue = \paramFunc(\paramValue)$ we have:
    \begin{equation}
        \pi(\xIt) = \pi \left( \rval{x} \mid \paramFuncValue \right)\cdot
        \label{eq:func_of_theta}
    \end{equation}
    The prior distribution corresponding to the central distribution for \set{X} in \Mt can be characterized
    by the following fixed point equation:
    \begin{equation}
        \bIaRoot \propto \left( \prod_i \frac{\abIab \rhoPrior}{\aIa} \right)^{\frac{1}{|\set{x}|}},
        \label{eq:fixed_point_cx}
    \end{equation}
    where
    \[
        \rhoPrior \propto \exp \left\{ \sumOverxe \xeIba \ln \bIxe \right\} \cdot
    \]

\end{theorem}

\begin{proof}
    First, we show that the set of distributions of each \Qvar is a convex set.
    Let $Q^*(\xIe) = \convexCond{\xIe}$, for some $Q_1, Q_2 \in \set{P}$, $0 \leq \alpha \leq 1$ and
    arbitrary $i, \xe$.

    \Mt is a parametric model.
    Therefore,
    \begin{align*}
        Q^*(\xIe) &= \alpha \int \xIebFull q_1(\beta \mid \rval {e}) d\beta + (1 - \alpha) \margDist[q_2]{\rval{x}}\\
        &= \int \pi(\xIt) (\convexCond[q]{\paramValue}) d\paramValue\cdot
    \end{align*}
    $\convexCond[q]{\paramValue}$ is a valid prior distribution and \Mt is a generic parametric model.
    That implies $Q^* \in \set{P}$.

    Thus, by Theorem~\ref{th:lambda} the central distribution of \set{X} satisfies
    \begin{equation}
        \Dparam[\central] = \lambda\cdot
        \label{eq:lambda_th_proof}
    \end{equation}
    Moreover, from Equation~\ref{eq:func_of_theta} and Bayes' theorem it follows
    \begin{align*}
        \lambda
        &= \Dparam[\central]\\
        &= \sumOverixe \xeIba \ln \frac{\xIeb}{\central(\xIe)} \\
        &= \sumOverixe \xeIba \ln \frac{\bIxe}{\cMass(\paramB \mid \rval{e})} \\
        &= \sum_i \ln \frac{\rhoPrior}{\bIa}\cdot
    \end{align*}
    Since each $\paramFunc$ is a function of $\param$, for $\paramFunc(\paramValue) = \paramFuncValue$
    we have $\cMass{\paramValue, \paramA, \paramB}=\cPrior$ $\cMass{\alpha , \paramA} = \cMass{\alpha}$.
    Therefore,
    \begin{equation*}
        \bIa = \frac{\aIa}{\abIab} \bIaRoot \cdot
    \end{equation*}
    Now we rewrite Equation~\ref{eq:lambda_th_proof}:
    \[
        |\set{X}| \ln \bIaRoot = -\lambda + \sum_i \ln \frac{\abIab \rhoPrior}{\aIa}
    \]
    Then:
    \begin{align*}
        \bIaRoot
        &= \left( e^{-\lambda} \prod_i \frac{\abIab \rhoPrior}{\aIa} \right)^{\frac{1}{|\set{x}|}} \\
        &= \left( e^{-\lambda} \prod_i \frac{\bIaRoot \rhoPrior}{\bIa} \right)^{\frac{1}{|\set{x}|}} \\
        &= \bIaRoot \left( \prod_i \frac{ \rhoPrior}{\bIa} \right)^{\frac{1}{|\set{x}|}} \\
        &= \bIaRoot \left(e^{-\lambda} \prod_i
        \frac{\rhoPrior[(\paramA, \paramB) \cMass(\paramA)]}{\cMass(\paramA, \paramB) \rhoPrior[(\paramA)] } \right)
        ^{\frac{1
        }{|\set{x}|}}
        \qedhere
    \end{align*}
\end{proof}

\bibliography{main}
\bibliographystyle{plain}